\section{Observations}\label{sec:observations}
\subsection{COSMOS2020}
We use the \cosmos Classic catalog and its released LePhare v2.1 redshift
probability distributions \citep{weaver2022,ilbert2006}. We select galaxies
with \texttt{flag\_combined}=0, \texttt{lp\_type}=0, HSC $i_{\rm ap}<25$ AB
in a $2''$ aperture, and native $i$-band signal-to-noise ratio $\geq5$.
We require finite CFHT $u$/HSC $grizy$ fluxes, positive flux uncertainties,
and a finite, nonnegative redshift PDF with nonzero integral. Negative fluxes
outside the $i$ band are retained. Of 170,003 objects passing the magnitude
and signal-to-noise cuts, \NParent\ have usable six-band measurements;
all have valid PDFs. Selection does not require spectroscopy.

We divide the field into $0.05^\circ$ sky cells and randomly assign cells
to training, validation, and test sets with probabilities 0.70, 0.15, and
0.15, using seed 20261024. We fix these assignments before magnitude
selection. The resulting sets contain \NTrain, \NValidation, and
\NPhotTest\ galaxies. No galaxy in a test cell is used for training.

\subsection{LSST-like Broadband Photometry}
The deep COSMOS observations provide both the multiband redshift labels and
the broadband measurements used to train our estimator. Using their native
broadband precision would test performance at COSMOS depth, rather than at
the noise levels expected for LSST. Photometric noise broadens color--redshift
degeneracies and changes the conditional distribution $p(z\given x)$.
We therefore add noise to the broadband inputs while retaining the redshift
information from the deep multiband observations. This lets us test how well
that information transfers to noisier broadband measurements and evaluate
all broadband methods on the same inputs.

We use native CFHT $u$ and HSC $grizy$ aperture measurements and add noise
to approximate LSST ten-year depths. The adopted $5\sigma$ point-source
depths are $(25.6,26.9,26.9,26.4,25.6,24.8)$ AB \citep{rubindepths}.
For each band, we set
\begin{align}
\sigma_{\rm tar}&=\frac{10^{(23.9-m_5)/2.5}}{5}\ \mu{\rm Jy},\\
\sigma_{\rm eff}&=\max(\sigma_{\rm deep},\sigma_{\rm tar}),\\
f'&=f_{\rm deep}+\epsilon,\quad
\epsilon\sim\mathcal N(0,\sigma_{\rm eff}^2-\sigma_{\rm deep}^2).
\label{eq:kernel}
\end{align}
We draw noise independently between bands and retain measurements already
noisier than the target depth. Selection precedes noise addition.
This model changes the noise level but retains the native filters.
It does not reproduce Rubin bandpasses, detection, blending, or correlated
measurement errors.

\subsection{Spectroscopic Redshifts}
We match Classic identifiers to the DR1.1 unique-object COSMOS spectroscopic
compilation \citep{khostovan2025}. We retain public galaxy measurements with
quality flags 3, 4, 13, or 14 and positive finite spectroscopic redshifts.
The parent match also requires a valid nonnegative catalog photo-$z$.
For duplicate observations, we use the compilation's preferred measurement.
We impose no photo-$z$/spec-$z$ agreement cut. The labels combine programs
including DESI, zCOSMOS, DEIMOS, C3R2, and LEGA-C.

The training, validation, and test sets contain 22,604, 4,803, and \NTest\
secure spectroscopic matches. We evaluate all methods on these same
\NTest\ test galaxies. Spectroscopy strongly favors brighter objects:
the median $i_{\rm ap}$ is 22.69 for the spectroscopic test sample and
24.17 for the full photometric test sample. Appendix~\ref{sec:selection}
shows their magnitude distributions.

\subsection{Benchmark Photo-zs}
We compare our method with three benchmarks:
\begin{itemize}
\item \textbf{Broadband flow.} A conditional flow trained on secure
spectroscopic redshifts. It uses the same architecture, six-band inputs,
preprocessing, noise augmentation, optimizer, and stopping rules as our
method. Only the training population and redshift labels differ.
We use this benchmark instead of SOMPZ because it directly compares
individual conditional redshift distributions within one model family.
SOMPZ groups galaxies in photometric space to calibrate ensemble redshift
distributions \citep{buchs2019,myles2021}. The flow avoids this discretization
and isolates differences from training data without introducing a different
estimator. We do not test whether it outperforms SOMPZ.

\item \textbf{Broadband EAZY.} Template fitting with EAZY-py 0.8.7
\citep{brammer2008}. We fit the saved noisy test measurements using its
12 tweaked FSPS templates, nonnegative template coefficients, 1\% systematic
flux error, and template-error amplitude 0.20. The redshift grid spans
$0.001<z<10$ with $\Delta z/(1+z)=0.005$. We correct fluxes and errors for
catalog Galactic extinction and apply the supplied
\texttt{prior\_R\_zmax7.dat} to HSC $r$, with an aperture-to-total correction.
This filter mapping is approximate. We fit no population-specific prior,
zero-point, or template corrections and provide no test redshifts to EAZY.

\item \textbf{Full-multiband LePhare.} The released COSMOS2020 redshift
distribution.
It supplies our training labels and uses more photometric information than
the broadband methods. Its comparison measures how much accuracy is retained
when transferring to six noisy bands.
\end{itemize}
Both template-fitting benchmarks have calibration histories that differ
from those of our flows.
